\documentclass{beamer}

\usepackage{beamerthemeSingapore}
\usepackage[brazil]{babel}
\usepackage[T1]{fontenc}
\usepackage[utf8]{inputenc}
\usepackage{graphicx}
\usepackage{url}
\usepackage{xcolor}
\usepackage{booktabs}
\usepackage{array}
\usepackage{ragged2e}
\usepackage{etoolbox}
\usepackage{amsmath}
\usepackage{minted}
\usepackage{tikz}
\usetikzlibrary{positioning, arrows.meta, calc}
% Grupo da Swiss Table, compartilhado com o roteiro.
% Grupo de 8 posições de uma Swiss Table e a busca por Davi, com etiqueta 89.
% Compartilhado entre o roteiro (figuras.tex) e os slides (hash.tex). Os
% valores das etiquetas são inventados para o exemplo.

\tikzset{
  swctrl/.style = {draw=black, fill=black!5, minimum width=1.25cm,
    minimum height=0.6cm, inner sep=0pt, font=\ttfamily\small, line width=0.6pt},
  swslot/.style = {draw=black, fill=white, minimum width=1.25cm,
    minimum height=0.75cm, inner sep=0pt, font=\ttfamily\small, line width=0.6pt},
  swrot/.style = {font=\sffamily\small, anchor=east},
  swidx/.style = {font=\sffamily\scriptsize, text=black!60},
  swnota/.style = {font=\sffamily\scriptsize, align=center},
  swseta/.style = {-{Stealth[length=2.2mm]}, line width=0.8pt}
}

% Marca de posição vazia no byte de controle: dois hífens, sem ligadura.
\newcommand{\vz}{\mbox{-}\mbox{-}}

\newcommand{\swissGrupo}{%
\begin{tikzpicture}
  \foreach \i/\c/\k in {0/12/Hugo, 1/\vz/, 2/89/Elisa, 3/45/Ana,
                        4/\vz/, 5/89/Davi, 6/3/Gabi, 7/\vz/} {
    \node [swidx] at (\i*1.25, 0.75) {\i};
    \node [swctrl] (c\i) at (\i*1.25, 0) {\c};
    \node [swslot] (s\i) at (\i*1.25, -1.0) {\k};
  }
  % Candidatas: etiqueta igual à procurada.
  \node [swctrl, fill=orange!35] at (c2) {89};
  \node [swctrl, fill=orange!35] at (c5) {89};
  \node [swslot, fill=red!12] at (s2) {Elisa};
  \node [swslot, fill=green!25] at (s5) {Davi};

  \node [swrot] at ($(c0.west)+(-0.15,0)$) {bytes de controle};
  \node [swrot] at ($(s0.west)+(-0.15,0)$) {posições};

  % A etiqueta procurada.
  \node [swctrl, fill=orange!35, label={[font=\sffamily\small]left:busca por \texttt{Davi}, etiqueta}]
    (q) at (2.5*1.25+1.25, 2.0) {89};
  \draw [swseta] (q.south) -- (c2.north east);
  \draw [swseta] (q.south) -- (c5.north west);
  \node [swnota, right=0.3cm of q] {compara com os 8 bytes\\de uma vez};

  \node [swnota, below=0.2cm of s2] {etiqueta igual,\\chave diferente};
  \node [swnota, below=0.2cm of s5] {chave\\encontrada};
\end{tikzpicture}}


\newcommand{\tcb}[1]{\textcolor{blue}{#1}}
\newcommand{\tcr}[1]{\textcolor{red}{#1}}
\newcommand{\tcg}[1]{\textcolor{green!50!black}{#1}}

\tikzset{
  slot/.style = {draw=black, minimum width=0.62cm, minimum height=0.55cm,
    inner sep=0pt, font=\sffamily\footnotesize, line width=0.6pt},
  idx/.style = {font=\sffamily\scriptsize, text=black!60},
  item/.style = {draw=black, fill=black!5, rounded corners=1pt,
    font=\sffamily\scriptsize, inner sep=2.5pt, line width=0.6pt},
  bit/.style = {draw=black, minimum width=0.42cm, minimum height=0.48cm,
    inner sep=0pt, font=\sffamily\scriptsize, line width=0.5pt},
  liga/.style = {-{Latex[length=1.6mm]}, line width=0.6pt}
}

\setminted{
  fontsize=\footnotesize,
  breaklines=true,
  breakanywhere=true,
  frame=lines,
  framesep=2mm,
  baselinestretch=1.05,
  tabsize=2
}

\apptocmd{\frame}{}{\justifying}{}

\setbeamertemplate{navigation symbols}{}

\title{10 - Tabelas de dispersão}
\author{Prof. Alex Kutzke}
\date{Outubro de 2026}

\begin{document}

% ==========================================
\begin{frame}
	\begin{center}
		\large{Setor de Educação Profissional e Tecnológica}\\
		Tecnologia em Análise e Desenvolvimento de Sistemas\\
		\vspace{1cm}
		\small{DS143 - Estruturas de Dados II}
	\end{center}
	\titlepage
\end{frame}

% ==========================================
\begin{frame}{Objetivos da aula}
	\footnotesize
	Ao final desta aula você deve ser capaz de:
	\begin{enumerate}
		\itemsep=0pt
		\item Explicar como a tabela de dispersão transforma a chave em posição, e
		      por que a busca não depende do número de chaves no caso médio;
		\item Avaliar uma função de dispersão quanto à distribuição das chaves, e
		      explicar por que $m$ costuma ser primo no método da divisão;
		\item Executar à mão a inserção com encadeamento e com sondagem linear,
		      sondagem quadrática e dispersão dupla;
		\item Relacionar o fator de carga ao custo médio da busca, e explicar
		      quando e por que a tabela é redimensionada;
		\item Explicar por que a remoção em endereçamento aberto exige lápide;
		\item Descrever o filtro de Bloom e o tipo de erro que ele admite;
		\item Escolher entre tabela de dispersão e árvore de busca a partir das
		      operações que a aplicação precisa.
	\end{enumerate}

	\vspace{0.1cm}
	Este conteúdo cai na Prova 2.
\end{frame}

% ==========================================
\begin{frame}{Roteiro}
	\tableofcontents
\end{frame}

% ==========================================
\section{Dispersão}

\begin{frame}[fragile]{Demonstração: o \texttt{map} de Go}
	\begin{minted}{go}
idade := map[string]int{
	"Ana": 19, "Bruno": 22, "Carla": 20, "Davi": 21,
	"Elisa": 23, "Fábio": 19, "Gabi": 24, "Hugo": 20,
}
for rodada := 1; rodada <= 3; rodada++ {
	for nome := range idade {
		fmt.Printf(" %s", nome)
	}
}
\end{minted}

	\begin{minted}{text}
$ go run mapa.go
  1: Fábio Gabi Hugo Ana Bruno Carla Davi Elisa
  2: Elisa Fábio Gabi Hugo Ana Bruno Carla Davi
  3: Carla Davi Elisa Fábio Gabi Hugo Ana Bruno
\end{minted}

	\vspace{0.1cm}
	Nem ordem alfabética, nem ordem de inserção, e cada percurso começa em uma
	chave diferente.
\end{frame}

\begin{frame}[fragile]{O tempo de uma busca}
	\begin{minted}{text}
Tempo médio de uma busca bem-sucedida, em nanossegundos:

         n | sequencial |  binária |     map
-----------+------------+----------+--------
      1000 |        429 |       82 |       6
     10000 |       1354 |       67 |       7
    100000 |      14014 |       92 |      13
   1000000 |     150637 |      158 |      29
\end{minted}

	\vspace{0.2cm}
	\begin{itemize}
		\item $n$ mil vezes maior: a sequencial fica 350 vezes mais lenta, a
		      binária, 2 vezes;
		\item o \texttt{map} é o mais rápido em todas as linhas;
		\item por que ele também cresce um pouco: no fim da aula.
	\end{itemize}
\end{frame}

\begin{frame}{Endereçamento direto}
	Chaves de 0 a 999: um vetor de 1000 posições, e a chave 427 fica em
	\texttt{v[427]}.

	\vspace{0.3cm}
	\begin{center}
		\begin{tikzpicture}
			\foreach \i/\t in {0/0, 1/1, 2/2, 3/\ldots, 4/427, 5/\ldots, 6/999} {
				\node [slot, minimum width=0.9cm] (d\i) at (\i*0.9, 0) {};
				\node [idx, below=0.05cm of d\i] {\t};
			}
			\node [slot, minimum width=0.9cm, fill=green!20] at (4*0.9, 0) {427};
		\end{tikzpicture}
	\end{center}

	\vspace{0.3cm}
	Busca, inserção e remoção em $O(1)$, sem comparação nenhuma.

	\vspace{0.4cm}
	Não serve quando o universo de chaves é grande:
	\begin{itemize}
		\item CPF tem 11 dígitos: 100 bilhões de posições para alguns milhares de
		      clientes;
		\item chaves de texto nem formam um intervalo de inteiros.
	\end{itemize}
\end{frame}

\begin{frame}{Tabela de dispersão}
	Uma \textbf{função de dispersão} (\emph{hash function}) leva cada chave a uma
	posição de um vetor de $m$ posições:
	\[
		h : \text{chaves} \to \{0, 1, \ldots, m-1\}
	\]

	\begin{center}
		\begin{tikzpicture}
			\foreach \i in {0,...,6} {
				\node [slot] (p\i) at (\i*0.62+3.5, 0) {};
				\node [idx, below=0.05cm of p\i] {\i};
			}
			\node [slot, fill=red!15] at (3*0.62+3.5, 0) {};
			\node [item] (a) at (4.6, 1.3) {\texttt{roma}};
			\node [item] (b) at (6.6, 1.3) {\texttt{rato}};
			\node [item] (c) at (3.0, 1.3) {\texttt{rei}};
			\draw [liga] (a.south) -- ([xshift=1.2mm]p3.north);
			\draw [liga] (b.south) -- (p6.north);
			\draw [liga, red] (c.south) -- ([xshift=-1.2mm]p3.north);
			\node [right=0.2cm of p6, font=\sffamily\scriptsize, text=red] {colisão em 3};
		\end{tikzpicture}
	\end{center}

	\vspace{0.1cm}
	\begin{itemize}
		\item $m$ é escolhido pelo tamanho do conjunto, e não do universo;
		\item duas chaves na mesma posição: \tcr{colisão};
		\item toda tabela de dispersão tem duas partes: a função e uma regra para
		      resolver colisões.
	\end{itemize}
\end{frame}

\begin{frame}{As colisões acontecem cedo}
	Chaves sorteadas em 365 posições: com \textbf{23 chaves}, a chance de alguma
	colisão já é de 50,7\%.

	\vspace{0.3cm}
	É o paradoxo do aniversário: em uma turma de 23 pessoas, a chance de duas
	fazerem aniversário no mesmo dia passa de metade.

	\vspace{0.4cm}
	\begin{center}
		\small
		\begin{tabular}{rr}
			\toprule
			\textbf{posições} & \textbf{chaves para 50\% de chance} \\
			\midrule
			365               & 23                                  \\
			1000              & 38                                  \\
			100003            & 373                                 \\
			\bottomrule
		\end{tabular}
	\end{center}

	\vspace{0.3cm}
	A regra de colisão não é um caso raro: ela trabalha o tempo todo.
\end{frame}

\begin{frame}{O que se exige da função de dispersão}
	\begin{block}{Determinística}
		A mesma chave cai sempre na mesma posição, senão a busca não acha o que a
		inserção guardou.
	\end{block}

	\begin{block}{Rápida}
		É calculada em toda operação.
	\end{block}

	\begin{block}{Uniforme}
		As chaves que a aplicação de fato usa se espalham pelas $m$ posições.
	\end{block}

	\vspace{0.2cm}
	A terceira depende das chaves, e não só da função: \texttt{go run funcoes.go}.
\end{frame}

\begin{frame}[fragile]{Método da divisão: $h(k) = k \bmod m$}
	\begin{minted}{text}
1. 1000 chaves múltiplas de 20, com h(k) = k % m

    m | posições usadas | chaves na posição mais cheia
------+-----------------+-----------------------------
  100 |               5 |                          200
   97 |              97 |                           11
\end{minted}

	\vspace{0.2cm}
	\begin{itemize}
		\item com $m = 100$, o resto de um múltiplo de 20 só pode ser 0, 20, 40, 60
		      ou 80;
		\item chaves que avançam de $p$ em $p$ ocupam $m / \text{mdc}(p, m)$
		      posições;
		\item com $m$ \textbf{primo}, o mdc é 1, e todas as posições são usadas.
	\end{itemize}

	\vspace{0.2cm}
	\footnotesize Padrões assim são comuns: identificadores de 10 em 10, preços
	terminados em 0 ou 9, endereços alinhados em 8.
\end{frame}

\begin{frame}[fragile]{Chaves de texto}
	\begin{columns}
		\begin{column}{0.45\textwidth}
			\begin{minted}[fontsize=\scriptsize]{text}
palavra | soma | polinomial
--------+------+-----------
amor    |   43 |         52
roma    |   43 |         62
mora    |   43 |         60
ramo    |   43 |          8
omar    |   43 |         77
\end{minted}
			\footnotesize A soma das letras ignora a posição: todo anagrama colide.
		\end{column}
		\begin{column}{0.53\textwidth}
			\begin{minted}[fontsize=\scriptsize]{go}
func polinomial(s string, m int) int {
	h := 0
	for i := 0; i < len(s); i++ {
		h = (31*h + int(s[i])) % m
	}
	return h
}
\end{minted}
			\footnotesize A palavra como um número na base 31, pelo método de Horner. É a
			base do \texttt{String.hashCode()} do Java.
		\end{column}
	\end{columns}

	\vspace{0.3cm}
	\[
		h(s) = \left(s_0 \cdot 31^{L-1} + s_1 \cdot 31^{L-2} + \cdots + s_{L-1}\right) \bmod m
	\]
\end{frame}

\begin{frame}[fragile]{Chaves parecidas entre si}
	\begin{minted}[fontsize=\scriptsize]{text}
3. 2000 matrículas, de GRR20260001 a GRR20262000, com m = 997

função          | posições usadas | chaves na posição mais cheia
----------------+-----------------+-----------------------------
soma das letras |              28 |                          150
polinomial      |             634 |                            6
FNV-1a          |             834 |                            7
sorteio         |             849 |                            7
\end{minted}

	\vspace{0.1cm}
	\begin{itemize}
		\item soma: só depende dos quatro dígitos finais, de 1 a 28;
		\item a polinomial espalha muito melhor, e ainda fica atrás das duas de
		      baixo;
		\item FNV-1a (pacote \texttt{hash/fnv}) mistura os bits com ou-exclusivo e
		      multiplicação;
		\item \texttt{sorteio} ignora a chave: é a referência de posições
		      independentes. Mesmo nela, 15\% das posições ficam vazias.
	\end{itemize}
\end{frame}

\begin{frame}[fragile]{A tabela usada nos exemplos}
	\begin{minted}[fontsize=\fontsize{7}{8}\selectfont]{go}
type Tabela struct {
	posicoes [][]Par
	n        int // número de pares guardados
}

func NovaTabela(m int) *Tabela {
	return &Tabela{posicoes: make([][]Par, m)}
}

func (t *Tabela) m() int {
	return len(t.posicoes)
}

func (t *Tabela) hash(chave string) int {
	h := 0
	for i := 0; i < len(chave); i++ {
		h = (31*h + int(chave[i])) % t.m()
	}
	return h
}
\end{minted}

	\vspace{-0.2cm}
	\footnotesize $m$ não é um campo: é o comprimento do vetor, e \texttt{m()} só o
	consulta. O \texttt{sondagem.go} tem a mesma organização, com o vetor de chaves
	no lugar do vetor de listas.
\end{frame}

% ==========================================
\section{Encadeamento}

\begin{frame}{Encadeamento}
	Cada posição guarda a \textbf{lista} das chaves que caíram nela.

	\vspace{0.3cm}
	\begin{center}
		\begin{tikzpicture}
			\foreach \i in {0,...,4} {
				\node [slot] (e\i) at (0, -\i*0.6) {};
				\node [idx, left=0.05cm of e\i] {\i};
			}
			\node [item, right=0.5cm of e0] (a1) {da};
			\node [item, right=0.5cm of e2] (b1) {a};
			\node [item, right=0.4cm of b1] (b2) {rei};
			\node [item, right=0.4cm of b2] (b3) {resto};
			\node [item, right=0.5cm of e3] (c1) {roeu};
			\node [item, right=0.4cm of c1] (c2) {roupa};
			\draw [liga] (e0.center) -- (a1);
			\draw [liga] (e2.center) -- (b1);
			\draw [liga] (b1) -- (b2);
			\draw [liga] (b2) -- (b3);
			\draw [liga] (e3.center) -- (c1);
			\draw [liga] (c1) -- (c2);
		\end{tikzpicture}
	\end{center}

	\vspace{0.3cm}
	\begin{itemize}
		\item inserção: calcula $h(k)$ e acrescenta à lista da posição;
		\item busca: calcula $h(k)$ e percorre \textbf{só} aquela lista;
		\item remoção: tira da lista.
	\end{itemize}
\end{frame}

\begin{frame}[fragile]{Busca no encadeamento}
	\begin{minted}[fontsize=\scriptsize]{go}
type Par struct {
	Chave string
	Valor int
}

func (t *Tabela) Busca(chave string) (int, bool) {
	for _, p := range t.posicoes[t.hash(chave)] {
		if p.Chave == chave {
			return p.Valor, true
		}
	}
	return 0, false
}
\end{minted}

	\footnotesize A lista é um \emph{slice}. Com lista encadeada, como no Cormen, o
	comportamento é o mesmo.
\end{frame}

\begin{frame}[fragile]{O programa: contar as palavras de uma frase}
	\begin{minted}{go}
t := NovaTabela(7)
for _, palavra := range strings.Fields(texto) {
	quantas, _ := t.Busca(palavra)
	t.Insere(palavra, quantas+1)
}
\end{minted}

	\begin{itemize}
		\item busca a contagem atual (zero, se a palavra é nova) e insere com a
		      contagem mais um. Chave já existente: \texttt{Insere} só troca o valor;
		\item a tabela começa com 7 posições. A frase tem 14 palavras distintas;
		\item a oitava, \texttt{de}, deixa 8 chaves em 7 posições, e a tabela
		      cresce:
	\end{itemize}

	\begin{minted}[fontsize=\scriptsize]{text}
inserção de "de"     n =  8 > m =  7, a tabela passa a ter 17 posições
\end{minted}
\end{frame}

\begin{frame}[fragile]{Contando as palavras de uma frase}
	\begin{columns}
		\begin{column}{0.5\textwidth}
			\begin{minted}[fontsize=\scriptsize]{text}
  0:
  1: [da 1]
  2:
  3: [roma 1]
  4:
  5: [de 2]
  6: [rato 2]
  7:
  8: [rainha 1]
  9: [o 2]
 10: [raiva 1]
 11:
 12: [a 2] [rei 1] [resto 1]
 13: [roeu 2] [roupa 2]
 14:
 15: [do 2]
 16: [e 1]
\end{minted}
		\end{column}
		\begin{column}{0.48\textwidth}
			\small
			\texttt{go run encadeamento.go}

			\vspace{0.3cm}
			\emph{o rato roeu a roupa do rei de roma e a rainha de raiva roeu o resto
				da roupa do rato}

			\vspace{0.3cm}
			14 chaves, 17 posições.

			\vspace{0.2cm}
			Seis posições vazias e uma lista de três chaves, com a tabela menos que
			cheia.
		\end{column}
	\end{columns}
\end{frame}

\begin{frame}{Fator de carga e custo}
	\[
		\alpha = \frac{n}{m}
	\]

	\vspace{0.1cm}
	No encadeamento, $\alpha$ é o comprimento médio das listas, e pode passar de 1.

	\vspace{0.3cm}
	Com \textbf{dispersão uniforme simples} (Cormen):
	\begin{itemize}
		\item busca malsucedida: percorre uma lista inteira, $\alpha$ chaves em
		      média;
		\item busca bem-sucedida: cerca de metade de uma lista;
		\item as duas custam $\Theta(1 + \alpha)$: o 1 é calcular $h(k)$.
	\end{itemize}

	\vspace{0.3cm}
	\begin{block}{Consequência}
		Com $\alpha$ limitado por uma constante, a busca custa $O(1)$ no caso
		médio. O pior caso continua $\Theta(n)$: todas as chaves na mesma posição.
	\end{block}
\end{frame}

\begin{frame}[fragile]{Redimensionamento}
	\begin{minted}[fontsize=\scriptsize]{go}
func (t *Tabela) redimensiona(m int) {
	nova := NovaTabela(m)
	for _, lista := range t.posicoes {
		for _, p := range lista {
			i := nova.hash(p.Chave)
			nova.posicoes[i] = append(nova.posicoes[i], p)
			nova.n++
		}
	}
	*t = *nova
}
\end{minted}

	\begin{itemize}
		\item \texttt{encadeamento.go}: cresce quando $\alpha > 1$, para o primeiro
		      primo depois de $2m$;
		\item $h$ depende de $m$: toda chave muda de posição e é reinserida;
		\item um redimensionamento custa $\Theta(n)$.
	\end{itemize}
\end{frame}

\begin{frame}{Custo amortizado}
	Tabela com 8 posições no início, que dobra quando recebe uma chave a mais do
	que tem posições. Inserindo 64 chaves:

	\vspace{0.1cm}
	\begin{center}
		\small
		\begin{tabular}{lrr}
			\toprule
			\textbf{inserção} & \textbf{posições depois} & \textbf{reinserções} \\
			\midrule
			9\textsuperscript{a}                & 16                       & 8                    \\
			17\textsuperscript{a}               & 32                       & 16                   \\
			33\textsuperscript{a}               & 64                       & 32                   \\
			\midrule
			\textbf{total}    &                          & \textbf{56}          \\
			\bottomrule
		\end{tabular}
	\end{center}

	64 inserções e 56 reinserções: 120 operações, menos de 2 por chave. A soma das
	reinserções é sempre menor que $n$, porque cada parcela é o dobro da anterior.

	\begin{block}{Custo amortizado}
		Custo médio por operação sobre a sequência inteira. A inserção custa $O(1)$
		amortizado, ainda que algumas custem $\Theta(n)$.
	\end{block}

	\footnotesize O \texttt{append} de Go funciona do mesmo jeito: \emph{slice}
	cheio, aloca um vetor maior e copia os elementos.
\end{frame}

% ==========================================
\section{Endereçamento aberto}

\begin{frame}{Endereçamento aberto}
	As chaves ficam no próprio vetor, sem listas. Posição ocupada: seguir uma
	\textbf{sequência de sondagem}, que depende da chave:
	\[
		h(k, 0),\; h(k, 1),\; h(k, 2),\; \ldots
	\]

	\vspace{0.2cm}
	\begin{itemize}
		\item a inserção para na primeira posição livre;
		\item a busca segue a mesma sequência, e para ao achar a chave ou uma
		      posição \textbf{vazia};
		\item a posição vazia prova que a chave não está: se estivesse, teria sido
		      posta ali ou antes;
		\item uma chave por posição: $\alpha < 1$ sempre.
	\end{itemize}
\end{frame}

\begin{frame}[fragile]{Três regras de sondagem}
	\begin{center}
		\small
		\begin{tabular}{lll}
			\toprule
			\textbf{Regra}      & $h(k, i)$                           & \textbf{a partir de $h_1(k)$} \\
			\midrule
			sondagem linear     & $(h_1(k) + i) \bmod m$              & $+0, +1, +2, +3, \ldots$      \\
			sondagem quadrática & $(h_1(k) + i^2) \bmod m$            & $+0, +1, +4, +9, \ldots$      \\
			dispersão dupla     & $(h_1(k) + i \cdot h_2(k)) \bmod m$ & passo $h_2(k)$                \\
			\bottomrule
		\end{tabular}
	\end{center}

	\begin{minted}[fontsize=\scriptsize]{go}
func (t *Tabela) posicao(k, i int) int {
	m := t.m()
	h1 := k % m
	switch t.regra {
	case Linear:
		return (h1 + i) % m
	case Quadratica:
		return (h1 + i*i) % m
	default: // Dupla
		h2 := 1 + k%(m-1)
		return (h1 + i*h2) % m
	}
}
\end{minted}
\end{frame}

\begin{frame}{Dispersão dupla: o passo $h_2(k)$}
	No \texttt{sondagem.go}:
	\[
		h_2(k) = 1 + (k \bmod (m-1))
	\]

	\vspace{0.1cm}
	$k \bmod (m-1)$ vai de 0 a $m-2$; somando 1, o passo fica entre 1 e $m-1$, e
	nunca é zero.

	\vspace{0.3cm}
	\begin{center}
		\small
		\begin{tabular}{cccl}
			\toprule
			\textbf{chave} & $h_1$ & $h_2$                 & \textbf{sequência de sondagem} \\
			\midrule
			16             & 5     & $1 + 16 \bmod 10 = 7$ & 5, 1, 8, 4, \ldots             \\
			27             & 5     & $1 + 27 \bmod 10 = 8$ & 5, 2, 10, 7, \ldots            \\
			\bottomrule
		\end{tabular}
	\end{center}

	\vspace{0.3cm}
	Com $m = 11$: o mesmo $h_1$ e sequências diferentes. Cada posição é a anterior
	mais o passo, módulo 11: para o 16, $5 + 7 = 12 \to 1$, $1 + 7 = 8$,
	$8 + 7 = 15 \to 4$.
\end{frame}

\begin{frame}[fragile]{Sondagem linear}
	\begin{minted}[fontsize=\scriptsize]{text}
$ go run sondagem.go
Sondagem linear, m = 11, h1(k) = k % 11
  insere 22: h1 =  0, posições examinadas: 1
  insere 33: h1 =  0, posições examinadas: 2
  insere  5: h1 =  5, posições examinadas: 1
  insere 16: h1 =  5, posições examinadas: 2
  insere 27: h1 =  5, posições examinadas: 3
  insere 38: h1 =  5, posições examinadas: 4
  insere 44: h1 =  0, posições examinadas: 3
      0   1   2   3   4   5   6   7   8   9  10
     22  33  44   .   .   5  16  27  38   .   .
\end{minted}

	\vspace{0.1cm}
	\small
	O 5, o 16, o 27 e o 38 têm $h_1 = 5$, e cada um examina uma posição a mais
	que o anterior.
\end{frame}

\begin{frame}{Agrupamento primário}
	\begin{center}
		\begin{tikzpicture}
			\foreach \i/\c in {0/22, 1/33, 2/44, 3/, 4/, 5/5, 6/16, 7/27, 8/38, 9/, 10/} {
				\node [slot] (q\i) at (\i*0.62, 0) {\c};
				\node [idx, below=0.05cm of q\i] {\i};
			}
			\foreach \i in {5,...,8} {
				\node [slot, fill=orange!30] at (\i*0.62, 0) {};
			}
			\foreach \i/\c in {5/5, 6/16, 7/27, 8/38} {
				\node [font=\sffamily\footnotesize] at (\i*0.62, 0) {\c};
			}
			\draw [liga, red] (q6.north) to [bend left=50] (q9.north);
			\node [font=\sffamily\scriptsize, text=red, above=0.45cm of q7] {$h_1 = 6$ vai parar no 9};
		\end{tikzpicture}
	\end{center}

	\vspace{0.3cm}
	\begin{itemize}
		\item as posições 5 a 8 formam um bloco contínuo de posições ocupadas;
		\item toda chave cujo $h_1$ caia no bloco vai parar no fim dele, e o
		      aumenta;
		\item bloco grande recebe mais chaves, e cresce mais depressa que bloco
		      pequeno.
	\end{itemize}

	\vspace{0.3cm}
	É o \textbf{agrupamento primário}, o defeito da sondagem linear.
\end{frame}

\begin{frame}[fragile]{Quadrática e dupla, mesmas chaves}
	\begin{minted}[fontsize=\scriptsize]{text}
Sondagem quadrática, m = 11, h1(k) = k % 11
      0   1   2   3   4   5   6   7   8   9  10
     22  33   .  38  44   5  16   .   .  27   .

Sondagem dupla, m = 11, h1(k) = k % 11
      0   1   2   3   4   5   6   7   8   9  10
     22  16  27  38  33   5   .   .   .   .  44
\end{minted}

	\begin{itemize}
		\item \textbf{quadrática}: o bloco não se forma, mas chaves com o mesmo
		      $h_1$ seguem a mesma sequência (agrupamento secundário). Com $m$ primo,
		      só há garantia de posição livre com $\alpha \le 1/2$;
		\item \textbf{dupla}: o passo depende da chave. O 16 anda de 7 em 7, o 27
		      de 8 em 8, o 38 de 9 em 9;
		\item $h_2(k)$ nunca pode valer 0, e precisa ser primo com $m$. Com $m$
		      primo e $1 \le h_2(k) \le m-1$, as duas condições valem.
	\end{itemize}
\end{frame}

\begin{frame}[fragile]{Custo em função do fator de carga}
	\begin{minted}[fontsize=\fontsize{6.5}{8}\selectfont, breaklines=false]{text}
Busca bem-sucedida
 alfa | linear | quadrática |  dupla | teoria linear | teoria uniforme
------+--------+------------+--------+---------------+----------------
 0.50 |   1.50 |       1.43 |   1.39 |          1.50 |            1.39
 0.75 |   2.50 |       1.96 |   1.85 |          2.50 |            1.85
 0.90 |   5.45 |       2.79 |   2.55 |          5.50 |            2.56
 0.95 |  10.23 |       3.49 |   3.16 |         10.50 |            3.15

Busca malsucedida
 alfa | linear | quadrática |  dupla | teoria linear | teoria uniforme
------+--------+------------+--------+---------------+----------------
 0.50 |   2.49 |       2.14 |   2.01 |          2.50 |            2.00
 0.75 |   8.62 |       4.47 |   4.01 |          8.50 |            4.00
 0.90 |  50.46 |      11.28 |   9.99 |         50.50 |           10.00
 0.95 | 203.56 |      22.97 |  19.87 |        200.50 |           20.00
\end{minted}

	\footnotesize Posições examinadas por busca, em média. $m = 100003$, chaves
	aleatórias. A linear examina mais posições e continua em uso: as posições
	vizinhas são lidas da memória em bloco, pelo mesmo motivo da página da árvore
	B, e na quadrática e na dupla cada posição pode exigir uma nova leitura.
\end{frame}

\begin{frame}{As fórmulas}
	\begin{block}{Sondagem linear (Knuth)}
		\[
			\text{bem-sucedida} \approx \frac{1}{2}\left(1 + \frac{1}{1-\alpha}\right)
			\qquad
			\text{malsucedida} \approx \frac{1}{2}\left(1 + \frac{1}{(1-\alpha)^2}\right)
		\]
	\end{block}

	\begin{block}{Dispersão uniforme, o caso ideal (Cormen, 11.6 e 11.8)}
		\[
			\text{bem-sucedida} \approx \frac{1}{\alpha} \ln \frac{1}{1-\alpha}
			\qquad
			\text{malsucedida} \approx \frac{1}{1-\alpha}
		\]
	\end{block}

	\vspace{0.2cm}
	\begin{itemize}
		\item a dupla fica sobre a curva ideal; a quadrática, um pouco acima;
		\item o $(1-\alpha)^2$ explica as 200 posições com a tabela 95\% cheia;
		\item crescer em: $1/2$ (Sedgewick, linear), $2/3$ (\texttt{dict} do
		      Python), $7/8$ (\texttt{map} de Go).
	\end{itemize}

\end{frame}

\begin{frame}{Remoção: esvaziar a posição não funciona}
	\begin{center}
		\begin{tikzpicture}
			\foreach \i/\c in {0/22, 1/33, 2/44, 3/, 4/, 5/5, 6/, 7/27, 8/38, 9/, 10/} {
				\node [slot] (r\i) at (\i*0.62, 0) {\c};
				\node [idx, below=0.05cm of r\i] {\i};
			}
			\node [slot, fill=red!15] at (6*0.62, 0) {};
			\draw [liga] (r5.north) to [bend left=50] (r6.north);
			\node [font=\sffamily\scriptsize, text=red, above=0.5cm of r6] {para aqui};
		\end{tikzpicture}
	\end{center}

	\vspace{0.3cm}
	O 38 foi inserido na posição 8 depois de passar por 5, 6 e 7. Se a remoção do
	16 esvaziar a posição 6, a busca pelo 38 para nela e responde que o 38
	\tcr{não está} na tabela.

	\vspace{0.4cm}
	\begin{block}{Lápide (\emph{tombstone})}
		Um terceiro estado, além de vazia e ocupada: a busca passa por cima, e a
		inserção pode reaproveitar a posição.
	\end{block}
\end{frame}

\begin{frame}[fragile]{Remoção com lápide}
	\begin{minted}{text}
Remoção com lápide, na tabela da sondagem linear:
  Remove(16) = true
  Busca(38) = true, posições examinadas: 4
      0   1   2   3   4   5   6   7   8   9  10
     22  33  44   .   .   5   X  27  38   .   .
\end{minted}

	\vspace{0.3cm}
	\begin{itemize}
		\item as lápides não contam em $n$, mas alongam as buscas como se fossem
		      chaves;
		\item muitas remoções acumulam lápides, e a tabela precisa ser reconstruída
		      de tempos em tempos, só com as chaves vivas.
	\end{itemize}
\end{frame}

\begin{frame}[fragile]{Remoção sem lápide, na sondagem linear}
	Esvaziar a posição e trazer para ela as chaves seguintes que dependiam dela:

	\vspace{0.2cm}
	\begin{itemize}
		\item percorre as posições depois do buraco até achar uma vazia;
		\item chave com $h_1$ igual ou anterior ao buraco desce para ele, e o buraco
		      passa para a posição que ela deixou.
	\end{itemize}

	\begin{minted}{text}
Remove(16), sem lápide:
      0   1   2   3   4   5   6   7   8   9  10
     22  33  44   .   .   5  16  27  38   .   .   antes
     22  33  44   .   .   5  27  38   .   .   .   depois
\end{minted}

	\begin{itemize}
		\item o 27 ($h_1 = 5$) desce para a 6, o 38 ($h_1 = 5$) para a 7, e a
		      posição 9 vazia encerra a remoção;
		\item a busca pelo 38 passa a examinar 3 posições;
		\item só na sondagem linear. Na quadrática e na dupla, a lápide é a solução
		      usual.
	\end{itemize}
\end{frame}

\begin{frame}{Encadeamento ou endereçamento aberto}
	\begin{center}
		\scriptsize
		\begin{tabular}{>{\raggedright\arraybackslash}p{2.2cm}
			>{\raggedright\arraybackslash}p{3.6cm}
			>{\raggedright\arraybackslash}p{3.6cm}}
			\toprule
			                 & \textbf{Encadeamento}                        & \textbf{Endereçamento aberto}                     \\
			\midrule
			Fator de carga   & pode passar de 1                             & sempre menor que 1, com folga                     \\
			$\alpha$ alto    & custo cresce devagar ($1 + \alpha$)          & custo dispara perto de 1                          \\
			Remoção          & tira da lista                                & exige lápide                                      \\
			Memória          & ponteiro ou \emph{slice} por posição, e nós  & só o vetor                                        \\
			Acesso à memória & segue ponteiros para outros endereços        & posições vizinhas, na linear                      \\
			Uso              & \texttt{HashMap} do Java                     & \texttt{dict} do Python, \texttt{map} de Go 1.24+ \\
			\bottomrule
		\end{tabular}
	\end{center}
\end{frame}

\begin{frame}{O pior caso provocado}
	Com a função polinomial de base 31:
	\[
		h(\texttt{Aa}) = 65 \cdot 31 + 97 = 2112 = 66 \cdot 31 + 66 = h(\texttt{BB})
	\]

	\vspace{0.1cm}
	\begin{itemize}
		\item \textbf{bloco}: cada um dos pares \texttt{Aa} e \texttt{BB}. Mesmo valor
		      e mesmo comprimento: trocar um pelo outro não muda o valor da palavra;
		\item \texttt{AaAa}, \texttt{AaBB}, \texttt{BBAa} e \texttt{BBBB} valem
		      todas 2031744;
		\item 10 blocos, cada um \texttt{Aa} ou \texttt{BB}: $2^{10} = 1024$ chaves
		      na mesma posição, e a tabela vira uma lista.
	\end{itemize}

	\vspace{0.3cm}
	\begin{block}{\emph{Hash flooding} (28C3, 2011)}
		Servidores web em PHP, Java, Python e outras linguagens derrubados por uma
		requisição com milhares de parâmetros que colidiam todos.
	\end{block}

	\vspace{0.2cm}
	Defesa: uma \textbf{semente} sorteada, misturada à função. O \texttt{map} de Go
	sorteia uma para cada \texttt{map} criado.
\end{frame}

% ==========================================
\section{Aplicações}

\begin{frame}{Tabela de dispersão e árvore de busca}
	A função espalha chaves vizinhas por posições distantes: a tabela perde a
	\textbf{ordem}.

	\vspace{0.3cm}
	\begin{center}
		\scriptsize
		\begin{tabular}{>{\raggedright\arraybackslash}p{2.8cm}
			>{\raggedright\arraybackslash}p{3.4cm}
			>{\raggedright\arraybackslash}p{3cm}}
			\toprule
			\textbf{Operação}        & \textbf{Tabela de dispersão}                       & \textbf{Árvore balanceada}  \\
			\midrule
			busca, inserção, remoção & $O(1)$ no caso médio, $\Theta(n)$ no pior          & $\Theta(\log n)$ no pior    \\
			mínimo e máximo          & $\Theta(n)$                                        & $\Theta(\log n)$            \\
			percurso em ordem        & ordenar antes: $\Theta(n \log n)$                  & $\Theta(n)$                 \\
			consulta por faixa       & percorre tudo                                      & desce ao início e segue     \\
			sucessor                 & percorre tudo                                      & $\Theta(\log n)$            \\
			\bottomrule
		\end{tabular}
	\end{center}

	\vspace{0.3cm}
	\small
	Frequência de palavras, sessões por identificador, duplicatas: só busca exata.
	Produtos entre 50 e 100 reais, próximo horário livre: precisa de ordem.
\end{frame}

\begin{frame}{Filtro de Bloom}
	Um vetor de $m$ bits, todos desligados no início, e $k$ funções de dispersão
	(Bloom, 1970).

	\vspace{0.3cm}
	\begin{itemize}
		\item inserir: calcula as $k$ posições e liga os $k$ bits;
		\item consultar: algum bit desligado, a chave \textbf{certamente não está};
		      todos ligados, \textbf{talvez esteja};
		\item o filtro não guarda as chaves, só os bits.
	\end{itemize}

	\vspace{0.4cm}
	\begin{block}{Um só tipo de erro}
		Uma chave nunca inserida pode achar todos os seus bits ligados por outras:
		\tcr{falso positivo}. Os bits de uma chave inserida nunca são desligados:
		falso negativo, nunca.
	\end{block}
\end{frame}

\begin{frame}{Exemplo: $m = 12$ bits, $k = 3$ funções}
	\begin{center}
		\begin{tikzpicture}
			\foreach \i/\b in {0/0, 1/1, 2/0, 3/1, 4/0, 5/0, 6/1, 7/0, 8/1, 9/0, 10/1, 11/0} {
				\node [bit] (b\i) at (\i*0.6, 0) {\b};
				\node [idx] at (\i*0.6, -0.45) {\i};
			}
			\node [item] (x) at (1.2, 1.4) {\texttt{casa}};
			\node [item] (y) at (5.1, 1.4) {\texttt{dado}};
			\foreach \i in {1,3,6} \draw [liga] (x.south) -- (b\i.north);
			\foreach \i in {3,8,10} \draw [liga] (y.south) -- (b\i.north);
			\node [item] (z) at (1.8, -1.6) {\texttt{gato}};
			\node [item] (w) at (5.1, -1.6) {\texttt{bola}};
			\foreach \i in {1,8,10} \draw [liga, orange!80!black] (z.north) -- ([yshift=-0.6cm]b\i.south);
			\foreach \i in {6,8,9} \draw [liga, red] (w.north) -- ([yshift=-0.6cm]b\i.south);
		\end{tikzpicture}
	\end{center}

	\vspace{0.1cm}
	\small
	\begin{itemize}
		\item inseridas: \texttt{casa} (1, 3, 6) e \texttt{dado} (3, 8, 10). O bit 3
		      foi ligado pelas duas, e o filtro não registra quem ligou;
		\item \tcr{\texttt{bola}} (6, 8, 9): bit 9 desligado, certamente não está;
		\item \textcolor{orange!80!black}{\texttt{gato}} (1, 8, 10): três bits
		      ligados, por \texttt{casa} e \texttt{dado}. Talvez: falso positivo.
	\end{itemize}

	\footnotesize Posições supostas para o exemplo.
\end{frame}

\begin{frame}[fragile]{Taxa de falsos positivos}
	\begin{columns}
		\begin{column}{0.42\textwidth}
			\begin{minted}[fontsize=\scriptsize]{text}
$ go run bloom.go
10 bits por chave

 k | medida | teoria
---+--------+-------
 1 |  9.18% |  9.52%
 2 |  3.20% |  3.29%
 3 |  1.73% |  1.74%
 4 |  1.16% |  1.18%
 5 |  0.93% |  0.94%
 6 |  0.82% |  0.84%
 7 |  0.81% |  0.82%
 8 |  0.82% |  0.85%
 9 |  0.87% |  0.91%
10 |  1.02% |  1.02%
\end{minted}
		\end{column}
		\begin{column}{0.55\textwidth}
			\[
				P \approx \left(1 - e^{-kn/m}\right)^k
			\]

			\small
			\begin{itemize}
				\item poucas funções: poucos bits conferidos;
				\item muitas: o vetor enche de bits ligados;
				\item melhor $k \approx \frac{m}{n}\ln 2 = 6{,}9$.
			\end{itemize}

			\vspace{0.2cm}
			\textbf{Bits por chave}: $m/n$. No exemplo, $12/2 = 6$; aqui, 10, com
			menos de 1\% de erro. Não depende do tamanho da chave, e a taxa de erro
			depende só dele e de $k$.

			\vspace{0.2cm}
			\footnotesize LevelDB e RocksDB podem guardar um filtro por arquivo em
			disco: a resposta \emph{certamente não está} economiza a leitura.
		\end{column}
	\end{columns}
\end{frame}

\begin{frame}[fragile]{De volta à demonstração: a ordem do \texttt{map}}
	Desde Go 1.24, o \texttt{map} é uma \textbf{Swiss Table}: endereçamento aberto.

	\vspace{0.2cm}
	\begin{itemize}
		\item uma semente sorteada por \texttt{map}: as posições mudam de uma
		      execução para outra;
		\item cada percurso sorteia a posição em que começa. As três linhas da
		      demonstração são a mesma sequência circular;
		\item a especificação diz que a ordem não é especificada, e o sorteio impede
		      que um programa dependa dela.
	\end{itemize}

	\vspace{0.2cm}
	Para percorrer em ordem, ordenar as chaves:
	\begin{minted}{go}
for _, nome := range slices.Sorted(maps.Keys(idade)) {
	fmt.Println(nome, idade[nome])
}
\end{minted}
\end{frame}

\begin{frame}{Swiss Table: um grupo de 8 posições}
	\begin{center}
		\resizebox{\textwidth}{!}{\swissGrupo}
	\end{center}

	\small
	\begin{itemize}
		\item grupo: 8 posições (par chave e valor) e 8 bytes de controle, um por
		      posição: vazia (\texttt{\vz}) ou 7 bits de $h(k)$;
		\item a parte alta de $h(k)$ escolhe o grupo; os 7 bits da parte baixa são
		      a \textbf{etiqueta};
		\item busca por \texttt{Davi}, etiqueta 89: uma comparação com os 8 bytes
		      acha as candidatas 2 e 5, e só nelas a chave inteira é comparada;
		\item \texttt{Elisa}: etiqueta igual por acaso, chance de 1 em 128.
	\end{itemize}
\end{frame}

\begin{frame}{De volta à demonstração: o tempo do \texttt{map}}
	\begin{itemize}
		\item chave fora do grupo e grupo com posição vazia: a chave não está;
		\item grupo cheio: passa a outro grupo, por sondagem quadrática;
		\item cresce com 7/8 das posições ocupadas. O fator de carga é alto porque
		      cada passo examina 8 posições de uma vez.
	\end{itemize}

	\vspace{0.4cm}
	\begin{block}{De 6 para 29 ns}
		O fator de carga fica sempre abaixo de 7/8, e a busca examina poucas
		posições em todas as linhas. Com um milhão de chaves, a tabela não cabe na
		memória \emph{cache}, e cada busca espera pela memória principal.
	\end{block}
\end{frame}

% ==========================================
\section{Fechamento}

\begin{frame}{Escolher entre as estruturas}
	\begin{center}
		\scriptsize
		\begin{tabular}{ll>{\raggedright\arraybackslash}p{4.6cm}}
			\toprule
			\textbf{Estrutura}           & \textbf{Busca}       & \textbf{Escolha quando}                   \\
			\midrule
			AVL                          & $\Theta(\log n)$     & busca domina, ordem importa               \\
			Rubro-negra                  & $\Theta(\log n)$     & uso geral com ordem, muitas alterações    \\
			Árvore B / B+                & $\Theta(\log_m n)$   & a estrutura vive em disco                 \\
			Trie                         & $O(\text{tamanho})$  & chaves são cadeias, consulta por prefixo  \\
			\tcb{Tabela de dispersão}    & \tcb{$O(1)$ médio}   & \tcb{só busca exata, sem ordem nem faixa} \\
			\tcb{Filtro de Bloom}        & \tcb{$O(k)$}         & \tcb{pertinência aproximada, pouca memória} \\
			\bottomrule
		\end{tabular}
	\end{center}

	\vspace{0.4cm}
	A pergunta que decide entre tabela de dispersão e árvore: a aplicação precisa
	da ordem entre as chaves?
\end{frame}

\begin{frame}{Para a próxima aula}
	\begin{itemize}
		\item exercícios 1 a 12 do roteiro da aula, sem entrega e sem nota. O
		      conteúdo cai na Prova 2;
		\item rodar \texttt{funcoes.go}, \texttt{encadeamento.go},
		      \texttt{sondagem.go} e \texttt{bloom.go}, alterando as chaves e os
		      tamanhos;
		\item leitura: Sedgewick, seção 3.4; Cormen, capítulo 11.
	\end{itemize}

	\vspace{0.5cm}
	Próxima aula: aula prática de tabelas de dispersão em laboratório, com a
	\textbf{Atividade Avaliativa 3}.
\end{frame}

\end{document}
